\chapter{Methodology}
\label{sec:methodology}

\section{Architecture}
\label{sec:method:architecture}
The encoder follows GNN+ \citep{luo2024gnnplus}; hyperparameters come from the literature
and from the tuning sweep of \ref{sec:method:experiment:hp}. One architecture is used for
unreduced, sparsified, partitioned and summarized graphs alike, with two exceptions stated
below: partitioning changes the pooling path (\ref{sec:method:architecture:partitioning}),
and the exact-compression track changes the input schema
(\ref{sec:method:architecture:exact}).

\begin{itemize}
    \item \textbf{Model Type:} Edge-Aware Graph Convolutional Network (GCN+) for AIGs
    \item \textbf{Input Representation \& Encodings}
    \begin{itemize}
        \item \textit{Node Features:} 4-dimensional vector of member type counts
        (Constant, Primary Input, AND Gate, Primary Output). On an unreduced graph every
        node has exactly one type, so this is a one-hot vector; on a summarized graph it
        counts the members of each super-node. The one-hot case is the size-1 special case
        of the count case, which is what allows one set of weights to ingest reduced and
        full graphs alike (\ref{sec:method:summary:schema}).
        \item \textit{Edge Attributes:} 2-dimensional vector of polarity counts (Standard
        connection, Inverted signal), one-hot on an unreduced graph and a count of the
        merged parallel edges on a summarized one.
        \item \textit{Positional Encoding:} Level-based, 32-dimensional space
        \item \textit{Pre-Encoder Integration:} 4D nodes linearly projected to 128D $\rightarrow$ concatenated with 32D PE $\rightarrow$ 160D vector $\rightarrow$ stabilized via \texttt{GraphNorm}
    \end{itemize}

    \item \textbf{Core Architecture Scale}
    \begin{itemize}
        \item Depth: 4 stacked message-passing layers
        \item Width/Hidden Dimension ($d_h$): 128 (consistent across layers)
    \end{itemize}

    \item \textbf{Edge-Aware Message Passing}
    \begin{itemize}
        \item Dynamic integration of edge attributes directly into message formulation
        \item Linear projection of edge features $\rightarrow$ fuse with neighbor features $\rightarrow$ LeakyReLU - from GNNPlus
        \item Symmetric degree normalization: Disabled in final configuration - Worked better GNNPlus found the same out...
    \end{itemize}

    \item \textbf{Layer Block Composition (Per Layer)}
    \begin{itemize}
        \item Activation \& Regularization: LeakyReLU + Dropout ($p = 0.15$)
        \item Residuals: Skip-connections added post-message passing
        \item Feed-Forward Network (FFN): Two-layer linear expansion ($128 \rightarrow 256 \rightarrow 128$)
        \item Normalization: \texttt{LayerNorm} applied before and after FFN
    \end{itemize}

    \item \textbf{Graph-Level Readout \& Aggregation}
    \begin{itemize}
        \item \textit{Jumping Knowledge (JK):} \texttt{"sum"} aggregation across all 4 layers (maintains $d_{out} = 128$)
        \item \textit{Global Pooling:} Mean pooling down to single graph-level embedding
        \item \textit{Regression Head Setup:}
        \begin{itemize}
            \item MLP architecture bridging graph embedding to final prediction
            \item Intermediate compression: Linear layer ($128 \rightarrow 64$) + \texttt{ReLU} + Dropout ($p = 0.3$, deliberately higher than the $0.15$ used inside the encoder blocks)
            \item Final projection: Linear layer ($64 \rightarrow 1$)
            \item Output bounding: Terminal \texttt{Sigmoid} activation constrains prediction to $[0, 1]$ range (mapping directly to the optimizability percentage)
        \end{itemize}
    \end{itemize}
\end{itemize}

\subsection{Partitioning Graphs}
\label{sec:method:architecture:partitioning}
Partitioning is the one reduction family that requires an architectural change. Graphs are
partitioned offline and the edges spanning partitions are removed, which leaves the node
set intact but disconnects the graph into $k$ components. Pooling directly to the graph
level would then average over components the encoder never allowed to communicate, so the
readout is made hierarchical:

\begin{itemize}
  \item GNN encoder (applied on intra-partition edges only)
  \item Pooling from node representations to partition representations
  \item Pooling from partition representations to a single graph representation
  \item Prediction head applied to the graph representation to produce the output $\hat{y}$
\end{itemize}

Feature normalization is also made partition-aware: the input \texttt{GraphNorm} normalizes
within each partition rather than within each graph, so statistics do not leak across a
boundary the message passing itself cannot cross.

\subsection{Sparsified Graphs}
\label{sec:method:architecture:sparsification}
Sparsification requires \emph{no} architectural change, and stating that is a
methodological claim rather than an omission: a sparsified AIG is still an AIG over the
same feature schema, only smaller, so the encoder of \ref{sec:method:architecture} ingests
it unmodified. What changes is the data pipeline.

Reductions are stored as boolean masks and applied when a graph is loaded, rather than
baked into the cached graphs, so that one cached corpus serves every sparsification setting.
The two kinds of mask behave differently downstream: edge-mask methods (random edge dropout,
spanning forest) leave node counts untouched, while node-mask methods (PageRank pruning,
AND-gate-only) do not. Because batches are assembled to a node budget
(\ref{sec:method:experiment:batching}), the node count after masking --- not before --- is
what the batch planner must use.

\subsection{Summarized Graphs}
\label{sec:method:architecture:summarization}
Summarized graphs are also fed to the unmodified encoder, which is possible only because
the merge preserves the input schema as a superset: node and edge features become member
counts rather than one-hot indicators, and a super-node containing a single node reproduces
its original feature row exactly (\ref{sec:method:summary:schema}). Positional encodings are
pooled by member minimum, so a super-node enters the circuit at the level of its earliest
member.

The one property this buys, and the reason it is worth the constraint, is that a full graph
requires \emph{no} preprocessing before a summary-trained model can be queried on it: a full
graph already \emph{is} a valid graph of size-1 super-nodes. This is what makes the
cross-state inference of RQ5 a direct test rather than a comparison across two input
formats.

\subsection{Exact Compression Track}
\label{sec:method:architecture:exact}
One summarization method --- colour refinement at count-cap $\infty$
(\ref{sec:method:summary:wl}) --- is exact rather than approximate, and realising that
guarantee requires a second, narrower encoder variant. This subsection states why the
general architecture above is not sufficient for it.

The obstacle is edge multiplicity. When $k$ nodes are merged, several parallel edges may
collapse into one super-edge, and the aggregated message must equal the sum of the $k$
messages it replaces. Carrying the multiplicity in \texttt{edge\_attr}, as the general
merge does, scales a value that is fed \emph{into} the message nonlinearity, and a
nonlinearity does not decompose over sums --- so the result is close but not equal. Applying
the multiplicity as an edge \emph{weight} \emph{after} the nonlinearity does give equality.

That in turn requires removing the edge attributes, since the exact track has no separate
edge encoder to fold polarity through. Polarity is instead moved onto the nodes: because
fan-in is fixed by node type in an AIG (zero for constants and primary inputs, two for AND
gates, one for primary outputs), the count of a node's inverted incoming edges together with
its type fully describes its polarity pattern for a graph-level target. What is given up is
knowing \emph{which} specific fan-in is inverted --- acceptable for this task, and scoped to
this track only; the general pipeline keeps polarity as a per-edge relation.

The consequences are worth stating explicitly, because they cut in both directions:

\begin{itemize}
    \item \textbf{The guarantee.} A model trained on graphs coarsened by exact colour
    refinement to depth $d$ equal to the number of message-passing layers produces
    \emph{identical} outputs to the same model on the uncoarsened graph
    \citep{bollen2023exact}. The accuracy cost of this compression is zero by construction,
    which makes it the anchor of the study and the positive control for RQ5
    (\ref{sec:method:experiment:design}).
    \item \textbf{The cost.} The schema is not the general one, so feature-space
    compatibility no longer holds for free. A full graph must be converted to the exact
    schema before an exact-track model can be queried on it. For the other methods
    cross-state inference needs no such step; here it does, and results must say so.
\end{itemize}

\subsection{Baselines for Optimizability Prediction}
\label{sec:method:architecture:baselines}
RQ1 asks how well a GNN predicts optimizability, which is only meaningful against a
reference. Three tiers are used, chosen so that a poor result at one tier is
interpretable rather than merely disappointing.

\subsubsection{Tier 1: Trivial Predictors}
\label{sec:method:baselines:trivial}
Constant predictors (training-set mean and median) and a regressor on graph size alone.
These establish what $R^2$ means on this label: a model that only learns the target's
central tendency scores $R^2 \approx 0$, and a model that only learns ``bigger graphs
optimize differently'' is the weakest non-trivial hypothesis about the task. Any claim that
the GNN reads circuit \emph{structure} has to clear the size-only predictor.

\subsubsection{Tier 2: Standard Graph-Level Encoders}
\label{sec:method:baselines:standard}
GCN \citep{kipf2017gcn}, GraphSAGE \citep{hamilton2017graphsage} and GIN
\citep{xu2019gin} under the same pooling, head, and training configuration as the primary
model. These isolate the contribution of the edge-aware, positionally-encoded architecture
of \ref{sec:method:architecture} from the contribution of message passing as such.

\subsubsection{Tier 3: Published Circuit Models}
\label{sec:method:baselines:published}
Three models from the circuit representation learning literature, ported to this dataset:

\begin{itemize}
    \item \textbf{SynthNet} \citep{chowdhury2021openabc}, the QoR model released with
    OpenABC-D. It is the closest published analogue of this task: OpenABC-D's graph-level
    labels include percentage of nodes optimized, which is essentially the label used here.
    \item \textbf{HOGA} \citep{deng2024hoga}, hop-wise graph attention. Depth of field comes
    from precomputed multi-hop features rather than from stacked message passing, so the
    trunk never sees an edge index at training time --- a different scaling strategy against
    which input reduction can be contrasted.
    \item \textbf{DeepGate4} \citep{deepgate4_2025}, sparse-attention circuit representation
    learning. This is the architecture-scaling counterpart to this thesis's input-scaling
    approach, which makes it the most direct comparison available.
\end{itemize}

\paragraph{Comparability protocol.}
Only the data pipeline is shared: identical splits, caches and label definitions, since a
different split would make the comparison meaningless. Everything else --- model
hyperparameters, optimizer, loss, learning-rate schedule --- defaults to each paper's own
published values rather than to this project's configuration, because these are separate
papers with their own tuned setups. Where a paper publishes no value for a setting this
task requires, the deviation is recorded rather than silently absorbed; the two substantive
ones are below.

\paragraph{Batching deviations, and why they are deviations in the papers' direction.}
Both HOGA and DeepGate4 must be re-batched, for the same underlying reason: they were
published for per-node tasks, and a graph-level target requires every node of a graph in
one forward pass, so a graph cannot be split across batches.

HOGA upstream minibatches over \emph{nodes} (1024 nodes drawn from a single preprocessed
graph), not over graphs. Batching a fixed number of graphs here would make memory unbounded,
since graphs in this corpus average tens of thousands of nodes. Bounding nodes per batch is
therefore \emph{closer} to upstream's own batching unit than a fixed graph count would be.
DeepGate4 needs a much tighter budget for a different reason: its virtual edge set is far
denser than the circuit itself, and each attention layer materialises a message tensor over
those edges, so gradient checkpointing is what makes the baseline runnable at all.

In both cases a node budget alone leaves each optimizer step averaging far fewer graph-level
losses than the primary model, because one graph contributes one label regardless of its
size. Gradient accumulation restores a comparable effective batch. The budget and the
accumulation count are therefore a \emph{pair}: their product is the effective batch, and
changing one alone silently changes the optimization regime the pairing exists to hold
constant.

% Full derivations, measured virtual-edge counts, activation-memory estimates and the exact
% list of which hyperparameters are published vs assumed are in src/train_baseline.py's
% module docstring and src/baselines/*/regressor.py. Condense, do not re-derive.

\paragraph{A baseline that fails is a result, and needs care.}
The first full SynthNet run collapsed to a constant prediction --- flat validation loss
across epochs with per-step metrics still moving, which is the signature of a model
predicting the target mean. This is not obviously a porting error: OpenABC-D's own paper
reports negative $R^2$ on the split variant used here (unseen design, as opposed to unseen
recipe within a seen design), so the failure mode is documented upstream.

The comparison must nevertheless be made carefully, and this is a trap worth naming in the
text. Upstream z-scores its targets \emph{per design}, which removes between-design variance
from the denominator of $R^2$; upstream's $R^2$ therefore measures within-design explanatory
power while the one reported here is dominated by between-design variance. The two are not
the same quantity and their values must not be placed side by side. Only the sign transfers.

% Detail and the diagnostic script live in src/baselines/openabc_synthnet/DIAGNOSIS.md.
% TODO: run src/diagnose_synthnet_baseline.py against the checkpoint and quote the measured
% per-split target mean and std rather than the values inferred from the metrics.


\section{Graph Reduction}
\label{sec:method:reduction}
This section outlines the diverse reduction methods integrated and tested in this framework for compressing AIGs. The primary objective of applying these reductions is to mitigate severe memory bottlenecks (such as GPU OOM errors) and accelerate training, all while attempting to preserve the structural features necessary for accurate optimizability prediction.

\subsection{Partitioning}
\label{sec:method:reduction:partitioning}
Partitioning methods divide massive AIGs into smaller, distinct subgraphs to manage memory constraints. This process retains all original nodes, only the edges that span between different partitions are cut. The number of partitions, $k$, is determined dynamically for each graph by dividing the total number of nodes by a target partition size, and then clamping the result to fall within a predefined minimum and maximum range.

\subsubsection{Random Hashing}
\label{sec:method:partition:random}
This method uniformly assigns nodes to partitions at random, serving as a naive baseline for evaluating the efficiency and performance of more complex partitioning algorithms.

\subsubsection{METIS}
\label{sec:method:partition:metis}
Implemented via PyTorch Geometric \citep{fey2019pyg}, this approach uses the classic METIS
algorithm \citep{karypis1997metis} to create balanced graph partitions while explicitly
minimizing the overall number of edge cuts. Its objective is topology-blind: all cut edges
count equally, regardless of what they connect. That is precisely the assumption the
domain-informed variant below modifies.

\subsubsection{Naive Level Slicing}
\label{sec:method:partition:levelslice}
This technique partitions the graph strictly by topological depth (e.g., grouping top-level and bottom-level nodes together). While naive, it helps preserve the circuit's directed logical flow better than random assignment.

\subsubsection{Span-Weighted METIS (Domain-Informed)}
\label{sec:method:partition:spanmetis}
This is a targeted, domain-aware adaptation of METIS. In AIGs, cutting edges that span across multiple topological levels breaks the longest causal chains. To prevent this, edges are assigned weights proportional to their span length (i.e., a large weight to long edges spanning multiple levels, and a small weight to short edges). METIS is then forced to minimize the \textit{weighted} cut, actively discouraging the severing of logic paths.


\subsection{Sparsification}
\label{sec:method:reduction:sparsification}
Sparsification reduces graph density by selectively removing edges or nodes while attempting
to preserve overall structural properties. Two of the four methods below remove edges and
leave the node set intact; the other two remove nodes and their incident edges. That
difference determines which retention figure is meaningful for each method, and it is why
\ref{sec:prelim:reduction:measuring} insists on reporting node and edge retention separately
rather than as one compression number.

\subsubsection{Random Edge Dropout}
\label{sec:method:sparse:dropout}
This naive baseline uniformly drops a fixed random fraction of edges across the graph,
acting as a control for structural degradation. The drop rate is the one freely tunable
compression knob in this family, which is what makes it usable as a matched-compression
control against the parameter-free methods below.
% Configured rate: 0.3 (config.SPARSIFICATION_RANDOM_DROPOUT_RATE), giving a measured 69.7\%
% edge retention.
% ACTION REQUIRED (see \ref{sec:results:rq4:pairings}): to pair this against spanning forest
% at matched compression the rate must be raised to approximately 0.419, matching spanning
% forest's measured 58.1\% edge retention. At 0.3 the two are not comparable. Changing it
% invalidates the existing random_edge_dropout runs, so decide before the next sweep.

\subsubsection{And-Gate Only}
\label{sec:method:sparse:andgate}
This domain-specific method reduces the graph by dropping all Primary Input (PI) and Primary
Output (PO) nodes, preserving only the internal AND logic gates. The rationale is that the
PI/PO boundary is fixed by the circuit's interface and cannot be restructured by synthesis,
so the AND network carries the optimizable structure. The method is parameter-free: its
compression is whatever fraction of the graph the interface happens to occupy (measured
82.1\% node and 73.3\% edge retention), not a value that can be dialled.

\subsubsection{Random Spanning Forest}
\label{sec:method:sparse:forest}
This method extracts a sparse connectivity backbone by computing a minimum spanning forest
over randomised edge weights, which eliminates reconvergent paths while keeping the graph's
connectivity intact. It is computed on the undirected projection of the AIG, so edge
direction plays no part in which edges survive --- a real loss of information on a DAG, and
one reason to expect this method to damage the predictive signal.

It is also parameter-free: a spanning forest has exactly $|V| - c$ edges for $c$ components,
so its compression (measured 58.1\% edge retention) is a property of the graph rather than a
setting.

This method replaced a stretch-bounded spanner, which was implemented and abandoned: spanners
exploit dense cyclic redundancy, and AIGs have little of it, so the spanner retained nearly
every edge. That negative result is worth reporting --- it is a small, concrete instance of
the thesis's general claim that generic graph machinery does not transfer to AIGs unexamined.
% The unused config.SPARSIFICATION_SPANNER_STRETCH constant is a leftover of that attempt.

\subsubsection{PageRank Node Pruning}
\label{sec:method:sparse:pagerank}
This centrality-based approach computes the PageRank score of every node and retains only
the highest-scoring fraction, implicitly dropping all edges incident to the removed
low-centrality nodes. The retention fraction is a knob (configured at 0.8, with damping
$\alpha = 0.85$), which is what allows this method to be matched against AND-gate-only:
0.8 was chosen to sit at that method's measured 82.1\% node retention, making the two a
directly comparable domain-informed / generic pair (\ref{sec:results:rq4:pairings}).


\subsection{Summarization \& Coarsening}
\label{sec:method:reduction:summarization}
These methods compress the graph by merging nodes into super-nodes. Unlike the other two
families, which only delete, summarization must decide what a merged node \emph{is} --- so
the shared rewrite that turns a clustering into a graph is described first
(\ref{sec:method:summary:mergemap}), and each method below then reduces to the single
question of which nodes it groups together.

The methods span a deliberate spectrum from fully domain-specific to domain-blind, so that
RQ4 has something to compare. \ref{tab:summarization_methods} lists the set.

\begin{table}[htbp]
\centering
\caption{Summarization methods. ``Run'' marks those given precompute and training runs;
the remainder are implemented and reported as literature context.}
\label{tab:summarization_methods}
\small
\begin{tabular}{llcc}
\toprule
\textbf{Method} & \textbf{Role} & \textbf{Compression} & \textbf{Run} \\
\midrule
Cone coarsening      & domain-specific       & band knob   & yes \\
MFFC skeleton        & domain-specific       & fixed       & yes \\
Colour refinement    & adapted, lossless     & fixed       & yes \\
ConvMatch            & GNN-aware SOTA        & target      & yes \\
Random within-type   & naive floor / matcher & exact       & yes \\
Spectral / local var.& domain-blind control  & target      & no  \\
LSH hashing          & cheap control         & bounded     & no  \\
\bottomrule
\end{tabular}
\end{table}

\subsubsection{The Merge Map and Super-Node Schema}
\label{sec:method:summary:mergemap}
Every method produces only a \emph{clustering}: an assignment of each node to a group. One
shared rewrite turns that clustering into a graph, so that methods differ in what they merge
and in nothing else. Given a clustering, the rewrite:

\begin{itemize}
    \item sets each super-node's features to the \emph{member type counts}
    $[\#\mathrm{const}, \#\mathrm{PI}, \#\mathrm{AND}, \#\mathrm{PO}]$;
    \item remaps every edge to its endpoints' clusters, discards the resulting self-loops,
    and sums the polarity counts of parallel super-edges so that a super-edge records how
    many normal and how many inverted connections it stands for;
    \item pools the positional encoding by member \emph{minimum}, so a super-node sits at
    the level of its earliest member;
    \item records the number of discarded intra-cluster edges, and the preserved primary
    input and output counts, as super-node and graph attributes.
\end{itemize}

\paragraph{Schema compatibility.}
\label{sec:method:summary:schema}
The count representation is a strict superset of the one-hot representation: a super-node
with one member reproduces its original feature row exactly. An unreduced AIG is therefore
already a valid input in the summarized schema, without conversion. This is the property
that makes RQ5 testable with a single set of weights, and it is a design constraint on every
method rather than a convenience --- a merge that produced features outside this space would
put summary and full graphs in different input spaces and make cross-state inference
undefined.

\paragraph{Boundary preservation.}
No method merges a primary input or primary output into a super-node with a different role.
The PI/PO boundary is the circuit's fixed interface, and dissolving it changes what the graph
represents rather than merely coarsening it.

% NOT YET DONE, and it is a cheap win: internal_edges, num_pis and num_pos are computed and
% attached to the graph but the encoder never reads them, and node_size is not attached at
% all (it is recoverable as x.sum(1)). The level PE is min-pooled only -- no max/mean/var.
% \citet{generale2022scaling}'s ablation found super-node CONTENT critical, so this is the
% obvious next experiment. Until it is wired into the model, do NOT write as if super-node
% content is being exploited: the current text above is accurate, an earlier draft was not.

\subsubsection{Cone Coarsening (Domain-Specific)}
\label{sec:method:summary:cone}
The AIG-native method, built on levels and reconvergence structure. It merges along two
independent axes:

\begin{itemize}
    \item \textbf{Depth axis --- fanout-free chain contraction.} Maximal chains of
    single-fanout gates contract into one super-node. This reduces circuit \emph{depth},
    which is what connects this method to hypothesis H1: contracting a chain shortens the
    path a signal must travel through the encoder.
    \item \textbf{Width axis --- level-band reconvergence merge.} Gates within a bounded
    band of topological levels that share a common immediate \emph{post}-dominator merge
    together. This compresses parallel width while leaving critical-path depth untouched, so
    the level-based positional encoding stays exact.
\end{itemize}

Two properties of this method were established empirically during implementation and are
reportable findings in their own right, because in both cases the natural specification
fails on real AIGs.

First, \textbf{the width axis must use post-dominators, not dominators}. Grouping gates by
common immediate dominator is the textbook phrasing, and it collapses on an AIG: a gate deep
in the circuit is reachable from the input frontier along many independent paths, so its
immediate dominator is the virtual source for approximately 99\% of AND gates measured.
Grouping on that reduces the axis to ``merge every gate on a level''. The immediate
post-dominator --- the gate at which a node's fanout cone reconverges --- is the meaningful
relation, and is also what ``reconvergence coarsening'' should have meant. The mirror
degeneracy then appears on the output side: any gate whose fanout reaches two primary
outputs without passing through a common gate post-dominates to the virtual sink, which on a
multi-output AIG is roughly 73\% of AND gates, and grouping those by level merges gates in
different connected components. Such gates are therefore excluded from merging entirely, so
the axis only ever groups gates with a genuine shared reconvergence point.

Second, \textbf{level bands are fixed windows, not $\pm k$ neighbourhoods}. ``Within $\pm 1$
level'' is not transitive and so does not define a partition at all; the implementation
assigns each node to the window $\lfloor \mathrm{level} / (b+1) \rfloor$ for band width $b$.
A consequence worth reporting is that widening the band does \emph{not} monotonically
increase compression, because a window boundary can fall between two otherwise-groupable
levels.

At band width zero the method is provably acyclicity-preserving: the width axis merges only
within a level and every netlist edge runs to a strictly higher level, so no edge lies inside
a group; the depth axis contracts a cluster into the single target of its only outgoing edge,
which cannot close a cycle in a DAG; and the depth axis runs on the width quotient, which is
still a DAG. Wider bands give up both this guarantee and the exactness of the level encoding.

% Production setting: max_chain_length = 4, level_band = 0.
% NOT YET MEASURED ON THE REAL CORPUS. On uniform-random DAGs both axes do essentially
% nothing (such graphs contain zero fanout-free AND->AND chains); on a generator that draws
% fanins from recently-created gates the depth axis contracts 331 -> 233 nodes. Do NOT quote
% either number: run one shard first (sbatch --array=32 src/shell/precompute_summarization.sh)
% and read _summary_stats_cone_shard000.json. If compression is weak on real AIGs, the
% contribution has to lean on retention at low compression instead.

\subsubsection{MFFC Skeleton Coarsening (Domain-Specific, Parameter-Free)}
\label{sec:method:summary:mffc}
The second domain-specific method contracts each \emph{maximal fanout-free cone} into one
super-node. Every AND gate whose fanout lies entirely inside one cluster joins that cluster,
iterated to a fixed point, so each gate ends up in the cone of the first gate at which all
its fanout paths reconverge. What survives is exactly the multi-fanout skeleton of the
netlist plus the fixed PI/PO interface.

The domain argument is the tightest in the family, and it is a direct answer to the
objection this method invites. MFFCs are indeed prime targets for synthesis rewriting ---
\texttt{refactor} collapses and re-expresses one MFFC at a time --- and one might therefore
expect merging them to destroy exactly the signal the model must read. The claim here is the
opposite, and it is falsifiable: intra-cone wiring is precisely what local rewriting is
\emph{free to change}, so it is structural noise with respect to the label, whereas the
sharing structure that \emph{constrains} what rewriting can achieve is what survives the
merge. If this method retains the label well, that is evidence the predictive signal lives
in the sharing structure rather than in local wiring; if it does not, the objection stands.

Two properties distinguish it from cone coarsening, with which it forms an ablation pair
rather than a rival:

\begin{itemize}
    \item \textbf{It is a strict generalization of the depth axis above.} Cone coarsening's
    chain contraction is the capped, chain-shaped special case; iterating absorption on the
    quotient additionally captures cones whose interior reconverges, which a chain rule
    cannot see.
    \item \textbf{It is parameter-free.} There is no ratio, band, or probe count to
    calibrate, so it cannot be accused of being the one method tuned in its own favour ---
    a real defence given that every other method here carries at least one uncalibrated knob.
    The price is that its compression is a fixed property of the circuit and cannot be
    dialled to match another method (\ref{sec:method:summary:random}).
\end{itemize}

Acyclicity is preserved by the same argument as the depth axis, applied per round: a cluster
is absorbed only into the single cluster receiving all of its outgoing edges. The fixed point
is unique and independent of visit order, so the method is deterministic.

% Absorption is capped at 64 rounds: a redundant staircase shape unlocks one absorption per
% round, which would be O(n) rounds; real netlists measure 2-3. Any prefix of the iteration
% is a valid partition, so the cap is safe rather than approximate.
% NOT YET MEASURED on the real corpus -- run sbatch --array=192 and check both retention and
% whether any graph hits the round cap.
% Variant B (annotating each cone with its NPN function class, i.e. ABC's own rewrite-library
% index) is designed but NOT built: it needs a model-schema change. Do not write it up as done.

\subsubsection{Graded Colour Refinement / Bisimulation (Adapted, Lossless)}
\label{sec:method:summary:wl}
The exact anchor of the study. Nodes are merged when they agree under $d$ rounds of colour
refinement, with $d$ set equal to the encoder's message-passing depth so that merged nodes
are exactly those the network cannot distinguish. A count-cap $c$ interpolates between two
named endpoints \citep{bollen2023exact}: at $c = \infty$ refinement distinguishes nodes by
the full multiset of neighbour colours and the merge is \emph{lossless} for a summing GNN;
at $c = 1$ it sees only the set of neighbour colours, which is bisimulation --- coarser,
cheaper, and lossy for a counting model.

Two AIG adaptations are required. Node colours are initialised from the four-way type rather
than uniformly, and \textbf{edge polarity is treated as a distinct relation} so that a
normal and an inverted connection never contribute to the same colour class ---
without this, logically distinct subgraphs merge. Refinement runs backward from the primary
outputs by default; direction is a knob, and it changes both how much compresses and which
structure survives.

At $c = \infty$ this method is orthogonal to the optimization being predicted by
construction: it removes only what the model provably cannot see, so it cannot erase the
label. That is what makes it the positive control for RQ5 rather than merely another
contender, and realising the guarantee end to end requires the schema of
\ref{sec:method:architecture:exact}.

The obvious risk is that structural hashing has already removed the easy equivalences
(\ref{sec:prelim:algorithms:primitives}), since every corpus graph is strashed. What this
method finds is therefore \emph{residual} redundancy beyond one hop, and how much of it
exists is itself a reportable statistic about AIGs.

% Production setting: depth = NUM_LAYERS (4), count_cap = None (exact), direction = backward.
% OPEN: the residual-redundancy probe (d = 1..4 over ~1-10k graphs) has not been run. Depth 1
% should find almost nothing precisely because of strash; that is a reportable result either
% way. This also subsumes k-SNAP/IO-summary as its d=1 case -- cite \citep{tian2008ksnap},
% do not run it as a separate method.

\subsubsection{Convolution Matching (General SOTA Bar)}
\label{sec:method:summary:convmatch}
The strongest general-purpose competitor: ConvMatch \citep{dickens2024convmatch} merges
nodes that are equivalent with respect to the \emph{graph convolution operation} itself
rather than with respect to a graph property, and reports retaining roughly 95\% of GNN
performance at 1\% of original size on node classification.

It is included as the bar the domain-specific methods must clear, and it is the right bar
precisely because it is GNN-aware but domain-blind: if an AIG-specific method beats it, the
domain claim is earned against a method optimising the same objective, not against a
strawman.

Its behaviour on a deep, polarity-carrying DAG is untested in the literature, and one
property found here explains why it may differ from the published setting. Its merge cost is
an unweighted sum over the convolution output, so merging two low-degree nodes scores better
than merging two genuine convolution twins of higher degree, regardless of similarity ---
verified directly on a constructed graph, where a pair of degree-two primary inputs scored
better than a matched twin pair. Convolution equivalence therefore decides \emph{between}
pairs of comparable degree, not across them. On a graph with the degree profile of an AIG
this biases merging toward the input frontier.

% Production setting: reduction_ratio = 0.5, sgc_depth = NUM_LAYERS, num_probes = 8, seed 42.
% num_probes replaces the reference's exact kNN over the SGC embedding. Raising it 2 -> 8
% measurably improves the paper's own objective but makes this the most expensive method in
% the family (~25 s and ~1.6 GB per graph at the corpus's largest size), so precompute must
% be sharded by graph size. Checked against the authors' reference implementation.

\subsubsection{Random Within-Type Merging (Naive Floor and Matching Instrument)}
\label{sec:method:summary:random}
Uniformly random nodes of the same type are merged until the node count reaches a target.
Merges are restricted to one node type so that the comparison isolates \emph{which} nodes
are merged rather than confounding it with whether the interface survives, matching the
boundary guarantee of every other method.

This serves two purposes, and the second is the important one.

As a \textbf{naive floor}, it answers the question a reader will ask first: if a
sophisticated method does not beat random merging at the same compression, its
sophistication bought nothing. Without it, three strong methods landing within noise of each
other cannot be told apart from a task that is indifferent to how nodes are merged. It is
also the exact counterpart of random edge dropout in the sparsification family, so both
halves of the study have a floor constructed the same way.

As a \textbf{matching instrument}, it is what makes RQ4 answerable at all. Compression is a
free parameter for only some methods: ConvMatch and the spectral method take a target ratio
directly, cone coarsening has a coarse integer band, and MFFC coarsening, colour refinement
and the hashing method have no compression knob whatsoever. There is therefore no general
way to make two \emph{real} methods meet at the same ratio. Random merging reaches any
target exactly by construction, so each method is instead compared against a random control
run at that method's own achieved compression. The paired design replaces a matching problem
that has no solution.

% ~10 lines on top of the shared merge rewrite; clustering cost is negligible. NOT YET BUILT
% -- it needs a METHODS entry appended at index 6 (or 7 alongside mffc) plus its own
% precompute and training runs per matched point. The runs, not the code, are the cost.

\subsubsection{Methods Implemented but Not Run}
\label{sec:method:summary:notrun}
Two further methods are implemented, tested and registered, but are reported as literature
context rather than given training runs. This is a scoping decision about compute, not a
statement that they failed.

\textbf{Spectral / local-variation coarsening} \citep{loukas2019localvariation} contracts
pairs scored to preserve the Laplacian spectrum. It was the natural domain-blind control,
but ConvMatch is the better general bar for this study: the claim under test is about
coarsening \emph{for GNN training}, so the comparison should be against the GNN-aware state
of the art rather than against a method optimising the graph spectrum, which nobody argues
is the right objective for graph-level regression. Its cost profile also makes it awkward:
local variation costs roughly two orders of magnitude more than heavy-edge matching at four
thousand nodes, so above a node cap it falls back to heavy-edge matching --- which means
that on a corpus of large graphs it is in practice ``heavy-edge matching with a spectral
rule on small graphs'', and must be described that way if it is reported at all.

\textbf{LSH-based hashing} \citep{ugc2024} merges nodes whose feature-and-connectivity
descriptors collide under locality-sensitive hashing, in linear time with no eigendecomposition
or iterative refinement. It was intended as the cheap naive tier, and measurement showed it
was never independent of colour refinement: at its calibrated setting it lands exactly on the
partition of distinct local descriptors, which is a hand-built one-round colour refinement.
Its compression is also bounded at both ends by properties of the descriptor rather than by
its bin-width knob, and the upper bound binds on every AIG measured --- so it cannot be
dialled to a target ratio and cannot participate in a matched-compression comparison. Both
of these are findings about the method and are reported as such
(\ref{sec:results:rq2:offline}), with its compression and retention presented on its own
curve rather than in the matched table.

The bound that makes this concrete is itself worth reporting as a dataset statistic: only
around 37\% of nodes in an AIG have a structurally distinct local descriptor. That is direct
evidence of the redundancy the entire summarization argument depends on.

\subsubsection{Functional Reduction as a Leakage Control}
\label{sec:method:summary:fraig}
All methods above are \emph{structural}: they merge nodes on connectivity and type, never on
Boolean function. That boundary is deliberate, and it can be probed rather than merely
asserted.

FRAIG-style functional reduction \citep{mishchenko2005fraig} merges functionally equivalent
nodes by SAT sweeping --- which is to say it performs part of the optimization being
predicted. Running it as a summarizer should visibly destroy or leak the optimizability
label, and doing so would demonstrate empirically where the boundary between compression
and optimization lies, rather than leaving it as an assumption. This is a negative control,
not a sixth method.

% NOT STARTED. If it is not run, say in \ref{sec:discussion:limitations:methods} that the
% structural/functional boundary is argued rather than measured.


\section{Data}
\label{sec:method:data}

\subsection{Source Circuits}
\label{sec:method:data:sources}
% TODO(provenance): AIG_DATASET_README.md states the corpus "combines two AIG sources"
% and carries an open item to finish downloading the original OpenABC-D. Before writing:
%   1. Name both sources and how many designs each contributes (the 55-design figure used
%      in \ref{sec:method:data:tier0} needs to be traced to them).
%   2. Record licensing for each -- required for the ethics section
%      (\ref{sec:discussion:ethics}) and for any dataset release.
%   3. Give the size and domain range they span (arithmetic, control, crypto, ...), since
%      the design-level split (\ref{sec:method:experiment:splitting}) assumes designs are
%      genuinely distinct circuits rather than variants of each other.
% Do not write this section from the node counts alone; the provenance is the point.

\subsection{Dataset Generation}
\label{sec:method:data:generation}

\subsubsection{Random Structural Transformation}
\label{sec:method:data:generation:transform}
% How the source designs are expanded into 231,055 structurally distinct base graphs.
% The methodological question to answer explicitly: why randomised transformation rather
% than plain replication? Replication would give identical graphs and therefore identical
% labels, teaching the model nothing; random transformation produces circuits that differ
% structurally while remaining functionally related, which is what makes a design-level
% split meaningful (\ref{sec:method:experiment:splitting}) and what the leakage argument
% there depends on.
% State the transformation actually applied and its parameters -- see data/creation/.

\subsubsection{Target Synthesis Algorithms}
\label{sec:method:data:generation:algorithms}
Four synthesis scripts were applied during dataset generation: \texttt{Orchestrate},
\texttt{Deepsyn}, \texttt{Syn4} and \texttt{C2RS}. They differ in how they search: some
apply a fixed sequence of rewriting, refactoring and balancing passes, while others explore
seeded or scored variations. Command templates and per-algorithm parameters are given in
\ref{sec:apx:first_appendix}.

Only \texttt{Orchestrate} is used for training and evaluation in this thesis
(\ref{sec:intro:scope:algorithm}); the remaining three exist on disk and are the most
direct extension available (\ref{sec:conclusion:futurework:algorithms}).
% TODO: one sentence each on what the four actually do, from AIG_DATASET_README.md and
% abc.rc. C2RS in particular is an alias expanding to a long balance/resub/rewrite chain.

\subsection{Tiered Dataset Structure}
\label{sec:method:data:tiers}
Over 3.9 million total unique AIGs.

\subsubsection{Tier 0: Base Graphs}
\label{sec:method:data:tier0}
This originates from 55 designs randomly transformed into 231,055 base graphs (Tier-0).

\subsubsection{Tier 1: Single-Algorithm Optimization}
\label{sec:method:data:tier1}
Applying the 4 distinct target synthesis algorithms (Orchestrate, Deepsyn, Syn4, C2RS) to these base graphs yields exactly 231,055 primary (Tier-1) samples per algorithm.

\subsubsection{Tier 2: Sequential Optimization}
\label{sec:method:data:tier2}
In Tier-2, each resulting graph from Tier 1 is then optimized by Orchestrate.

\subsection{Target Label: Optimizability}
\label{sec:method:data:label}

\subsubsection{Definition}
\label{sec:method:data:label:definition}
For an AIG $G$ and synthesis algorithm $S$, the optimizability target is the relative node
reduction
\[
    y(G) \;=\; 1 - \frac{|V(S(G))|}{|V(G)|} \;\in\; [0, 1],
\]
formulated as a regression task. The bound motivates the terminal sigmoid in the prediction
head (\ref{sec:method:architecture}): the model cannot predict outside the achievable range.

Note what this does and does not measure. It is a \emph{relative} quantity, so it is
comparable across circuits of different sizes --- which is what makes a single model over a
heterogeneous corpus meaningful --- but it is also therefore harder to predict than an
absolute count, since the model must infer both what will be removed and what the graph
started with. It is an area proxy and says nothing about delay or power
(\ref{sec:intro:scope:target}).

\subsubsection{Label Distribution}
\label{sec:method:data:label:distribution}
% BLOCKED: needs a histogram plus per-tier summary statistics (mean, std, quantiles).
% This is not cosmetic -- three later claims depend on the numbers:
%   1. R^2 is a ratio against the target's own variance, so a narrow label distribution
%      makes a respectable RMSE compatible with a near-zero or negative R^2. The baseline
%      diagnosis in \ref{sec:method:architecture:baselines} turns on exactly this.
%   2. If the distribution is concentrated, a constant predictor is a strong baseline and
%      the tier-1 comparison in \ref{sec:results:rq1:baselines} carries most of the weight.
%   3. Tiers differ by construction: a graph already optimized once (tier 2) has less left
%      to remove than a base graph (tier 0), so the tiers should have visibly different
%      label distributions. If they do not, that is itself a finding about the pipeline.
% Report per tier AND pooled; the pooled distribution is what the model actually sees.

\subsection{Graph Representation}
\label{sec:method:data:representation}

\subsubsection{Node and Edge Features}
\label{sec:method:data:representation:features}
PyTorch Geometric graphs utilizing 1-hot node types (Constants, PI, AND, PO), 1-hot encoded edge polarity (regular connection or inverted signal), and topological depth encodings.

\subsubsection{Positional Encoding Precomputation}
\label{sec:method:data:representation:pe}
Topological levels are computed once, offline, and stored on the graph; the model reads the
stored attribute rather than recomputing it. This matters for more than speed. Levels are
computed on the \emph{unreduced} graph and then carried through reduction, so a partitioned
or sparsified graph retains the depth information of the circuit it came from even where the
edges that establish that depth have been cut. A runtime encoding would instead describe the
reduced graph's own topology, which is a different --- and, for this task, less
informative --- quantity.

\subsection{Dataset Statistics}
\label{sec:method:data:statistics}
% BLOCKED on a stats pass over the corpus. Report, per tier and pooled: node and edge count
% distributions (min/median/max and a histogram), depth distribution, and the ratio of AND
% gates to interface nodes.
% Three things downstream need specific numbers from here:
%   - the graph-scale bound in \ref{sec:method:experiment:scale} is justified by the node
%     count distribution, and currently states the bound without showing the distribution
%     it was drawn from;
%   - the node-budget batch sizes in \ref{sec:method:experiment:batching} only make sense
%     against the mean and maximum graph size;
%   - the AND-gate fraction is what determines and_gate_only's compression
%     (\ref{sec:method:sparse:andgate}), so it explains a measured retention figure rather
%     than merely accompanying it.
% Worth adding as a structural statistic, since the summarization argument rests on it:
% the fraction of nodes with a structurally distinct local descriptor
% (\ref{sec:method:summary:notrun}) and, if the probe is run, the residual redundancy after
% strash at refinement depths 1-4 (\ref{sec:method:summary:wl}).

\subsection{Held-Out Hyperparameter Tuning Subset}
\label{sec:method:data:hpsubset}
A smaller, balanced (between Tier 0, 1 and 2) subset of 50,000 graphs was used exclusively for the initial hyperparameter tuning and then excluded from the rest.


\section{Experiment}
\label{sec:method:experiment}

\subsection{Research Design}
\label{sec:method:experiment:design}
This is a controlled comparison against a fixed reference point. The manipulated variable is
the reduction method applied to the training data; the encoder, dataset, splits, label,
optimizer and evaluation protocol are held constant, so that differences in outcome are
attributable to the reduction. The reference point is the full-graph model of RQ1, and every
reduced configuration is reported relative to it rather than in isolation.

Two hypotheses are stated in advance (\ref{sec:intro:rq3}, \ref{sec:intro:rq4}): H1, that
path-contracting summarization retains more than edge-removing sparsification at matched
compression, because contraction shortens propagation paths where deletion only lengthens
them; and H2, that the domain-informed adaptations tested here are too simple to win
outright, in which case the informative outcome is what the null result says about where the
predictive signal lives.

\paragraph{Positive control.}
The design includes one configuration whose outcome is known in advance. Colour refinement at
count-cap $\infty$ on the exact-compression track is provably lossless for this encoder
(\ref{sec:method:architecture:exact}), so a model trained on those coarsened graphs must
score essentially identically to the full-graph baseline, both on reduced graphs and when
queried on full ones. This is what makes a poor cross-state result interpretable: without it,
RQ5 cannot distinguish ``models trained on reduced graphs do not generalize'' from ``the
evaluation path is wrong''. Any deviation of this configuration from the baseline is a defect
to be found, not a finding to be reported.

\paragraph{What this design cannot establish.}
The comparison is within one encoder family, one synthesis algorithm and one bounded graph
scale. It can rank reduction methods against each other under those conditions; it cannot
establish that the ranking transfers to another architecture, another target algorithm, or
industrial-scale circuits. \ref{sec:discussion:limitations} treats each in turn.

\subsection{Experimental Setup}
\label{sec:method:experiment:setup}

\subsubsection{Data Splitting}
\label{sec:method:experiment:splitting}
An 80/10/10 split by volume, partitioned \emph{by design}: all graphs derived from a given
source circuit fall in the same split. This is the strong choice and it is deliberate.
Because the corpus is built by transforming and then repeatedly optimizing a small number of
source designs, a random split would place a base graph in training and its own optimized
descendant in test --- structurally near-identical circuits in different optimization states.
The model could then score well by recognising the design rather than by reading its
structure. Splitting by design removes that path entirely, at the cost of making the task
harder: the test set contains circuits the model has never seen in any state.

Two alternative strategies are implemented and available (split by optimization recipe, and
a plain random split). They are not used for the headline results, but they bound the leakage
argument: the gap between a design-level and a random split is a direct measurement of how
much of the score is design recognition. Reporting that gap once would strengthen the
validity discussion in \ref{sec:discussion:limitations:validity} considerably.

\subsubsection{Graph Scale Bounds}
\label{sec:method:experiment:scale}
The corpus is bounded at 366,040 nodes per AIG. This intermediate scale is what makes RQ1
possible: a single unreduced graph fits in accelerator memory, so the full-graph baseline
every other result is measured against can actually be trained. Larger circuits would remove
the reference point and leave only reduced configurations to compare with each other.

The bound is a scoping decision with a cost, and it is the limitation closest to the
motivation: the argument for this work invokes circuits of millions to billions of nodes, and
the evaluation does not reach them
(\ref{sec:discussion:limitations:scalability}).

\subsubsection{Dynamic Batching}
\label{sec:method:experiment:batching}
Graphs in this corpus vary in size by orders of magnitude, so a fixed batch size in graphs
gives wildly varying memory use and eventually fails on the largest circuits. Batches are
instead assembled to a maximum \emph{node} budget by greedy bin-packing, which bounds
activation memory regardless of which graphs are drawn. The batch plan is computed once and
cached, since replanning over a corpus this size is expensive.

Two budgets are used and they are deliberately distinct. Training uses the smaller of the
two; evaluation, which stores no activations or gradients, uses a substantially larger one.
Keeping them separate rather than raising a single shared value avoids invalidating the batch
plan the trained checkpoints were produced under.

\paragraph{A constraint on every efficiency number.}
The evaluation budget must be identical across all configurations, because peak memory and
throughput both depend on it. Any change requires re-running every configuration; a mixed set
of budgets makes the RQ2 comparison meaningless rather than merely noisy. This is stated here
because it constrains how the results in \ref{sec:results:rq2} may be read and combined.

One irreducible peak remains that no budget can lower: a graph larger than the budget still
forms a batch of its own, since a graph-level target cannot be split across batches. The
largest circuit in the corpus therefore sets a floor on peak memory for every configuration
that does not reduce node counts.

\subsubsection{Model Configuration}
\label{sec:method:experiment:model}
An Edge-Aware Graph Convolutional Network (GCN+) is trained to predict the optimizability of
a circuit when subjected to the \texttt{Orchestrate} synthesis script. The architecture is
described in \ref{sec:method:architecture} and is held fixed across all reduction
configurations; the baselines it is compared against are described in
\ref{sec:method:architecture:baselines}.

\subsubsection{Hyperparameter Tuning}
\label{sec:method:experiment:hp}
Hyperparameters were selected by an Optuna \citep{akiba2019optuna} sweep on the held-out
balanced subset of \ref{sec:method:data:hpsubset}, which is then excluded from all
subsequent training and evaluation. Tuning on a subset that later appeared in training or
test would let the reported scores absorb information from the tuning process; excluding it
entirely costs 50,000 graphs and removes the question.

Tuning was performed once, on the full-graph configuration, and the resulting values are held
fixed for every reduced configuration. This is the conservative choice for the comparison
being made: re-tuning per reduction method would confound ``this reduction retains more
signal'' with ``this reduction received more tuning effort''. It does mean reduced
configurations are evaluated at hyperparameters chosen for a different input distribution,
which if anything biases against them --- worth stating explicitly, since it bounds the
strength of any negative result about reduction.
% TODO: record the search space and the number of trials from src/shell/big_hp_tuning.sh.

\subsubsection{Training Configuration}
\label{sec:method:experiment:training}
Smooth L1 loss with $\beta = 0.01$, chosen well below the default because the target occupies
a narrow band within $[0,1]$ and the default threshold would place nearly every residual in
the quadratic regime, making the loss effectively squared error. Optimization uses AdamW with
learning rate $3 \times 10^{-4}$ and weight decay $3.96 \times 10^{-3}$, a linear warmup over
the first 10,000 steps, and \texttt{ReduceLROnPlateau} thereafter (factor 0.5, patience 2,
floor $10^{-6}$). Training stops early on validation loss with patience 3. Gradients are
clipped at 1.0.

Two settings affect the measurements rather than the model and must therefore be disclosed
with them: \texttt{torch.compile} is enabled, which changes every timing figure, and mixed
precision is used on GPU, which changes every memory figure. Neither is varied across
configurations, so comparisons remain valid; absolute numbers are only meaningful together
with these settings.

\subsubsection{Hardware \& Software Environment}
\label{sec:method:experiment:hardware}
All experiments ran on the Snellius national supercomputer: H100 80\,GB GPU nodes for
training and evaluation, and CPU-only nodes for the offline reduction precompute, which is
embarrassingly parallel over graphs and does not benefit from a GPU.

Two measured characteristics of the workload are worth recording, because they explain the
shape of the RQ2 results rather than merely describing the setup. First, memory scales close
to linearly with nodes and edges per batch --- the encoder's own parameters are negligible at
this width and depth --- which is why a node budget bounds memory effectively at all. Second,
the kernels are latency-bound on message-passing gather/scatter rather than
compute-bound: at full streaming-multiprocessor occupancy the floating-point pipelines and
memory bandwidth both sit far below saturation. The consequence for RQ2 is that a larger
batch amortises per-batch overhead but does not deliver a proportional speedup, and that
reductions which remove \emph{edges} should help throughput more than their node compression
alone would suggest.

% TODO: library versions (torch, PyG, Lightning) from pyproject.toml, and the total
% wall-clock/GPU-hour budget consumed -- the latter is also needed for
% \ref{sec:discussion:ethics}.

\subsection{Evaluation Metrics}
\label{sec:method:experiment:metrics}
Four groups of metrics are reported. Accuracy and ranking are reported together and never
separately, for a reason established in the literature: CTS-Bench \citep{cts_bench2026}
found that on a coarsened EDA graph the mean absolute error barely moved while $R^2$ fell
below zero. Absolute error alone can therefore look healthy while the model has lost all
discriminative power --- exactly the failure a reduction study is most at risk of missing.

\paragraph{Evaluation protocol.}
Each trained configuration is evaluated in two passes over the test set: once on graphs under
the same reduction the model was trained with (\emph{matched state}), and once on unreduced
graphs (\emph{full graph}). The matched pass answers RQ3 and the pair together answers RQ5,
since the difference between them isolates the cost of the structural shift from the cost of
the reduction itself.

\subsubsection{Predictive Accuracy}
\label{sec:method:experiment:metrics:accuracy}
Smooth L1 loss (the training objective, reported for comparability with training curves),
root mean squared error, and the coefficient of determination $R^2$. RMSE is on the label's
own scale and so is directly interpretable as ``percentage points of optimizability''; $R^2$
is a ratio against the target's variance and answers the different question of how much of
the variation the model explains, with zero meaning no better than predicting the mean.

$R^2$ is computed for validation and test only, never per training batch: it is a ratio
against the variance within whatever sample it is computed on, and a single batch's variance
is not a meaningful denominator.

\subsubsection{Ranking Quality}
\label{sec:method:experiment:metrics:ranking}
Spearman rank correlation between predicted and true optimizability across the test set.
Ranking is reported separately from error because it can survive when absolute accuracy
degrades, and because it is closer to how such a model would be used: a practitioner
deciding where to spend synthesis effort needs the ordering of circuits, not calibrated
values.

The scope of this metric should be read carefully. It ranks \emph{circuits} under one fixed
script, not \emph{scripts} for one circuit --- the latter is the script-selection task the
motivation invokes and this thesis does not attempt
(\ref{sec:intro:scope:prediction}).

\subsubsection{Efficiency \& Resource Metrics}
\label{sec:method:experiment:metrics:efficiency}
Peak allocated GPU memory (allocated rather than reserved, since reserved reflects allocator
behaviour rather than the model's demand), training step and epoch wall-clock time, inference
throughput in graphs per second, and GPU utilisation. All are recorded under the fixed
evaluation batch budget of \ref{sec:method:experiment:batching}; comparing across
configurations run at different budgets would be invalid.

\subsubsection{Reduction Quality Metrics}
\label{sec:method:experiment:metrics:reduction}
Node and edge retention ratios, reported separately (\ref{sec:prelim:reduction:measuring}),
and the offline wall-clock cost of computing the reduction. The offline cost matters
independently of the training saving: a reduction that costs more to compute than it saves
in training is not useful unless the artifact is reused across many runs, which is precisely
what the cached-reduction design of this pipeline enables. That trade-off should be stated
as an amortisation threshold --- how many training runs a reduction must serve before it pays
for itself --- rather than as a raw cost.

\paragraph{Effective receptive field.}
Hypothesis H1 claims that coarsening helps by contracting paths, and that claim needs a
measurement of its own. The metric is the mean size of the $k$-hop fanin cone over sampled
nodes, with $k$ set to the number of encoder layers, computed before and after reduction.
This is well defined for every method, unlike a DAG longest-path metric, which is only
meaningful for the methods that guarantee an acyclic quotient.

% NOT YET IMPLEMENTED. Until it is, H1 is asserted and not evidenced -- see
% \ref{sec:intro:rq3}. Either build it or downgrade H1 to a discussion point; do not report
% it as tested.

\subsubsection{Statistical Treatment}
\label{sec:method:experiment:metrics:stats}
Efficiency comparisons are made pairwise per graph against the full-graph baseline, with
bootstrap confidence intervals and a Wilcoxon signed-rank test, since the per-graph savings
are paired by construction and are not normally distributed.

Accuracy comparisons currently have no equivalent treatment, and this is a real weakness
rather than an omission of detail. Each configuration is trained once, so there is no
run-to-run variance against which to judge whether a difference is real --- which bears
directly on RQ4, whose expected effect is small. Either the RQ4 pairings are re-run across
several seeds, or the accuracy differences must be reported as point estimates and read with
corresponding caution (\ref{sec:discussion:limitations:reproducibility}).

\subsection{Experimental Phases}
\label{sec:method:experiment:phases}
The work proceeds in four phases, which do not map one-to-one onto the five research
questions: Phase~1 answers RQ1, Phase~2 answers RQ2, Phase~3 answers both RQ3 and RQ4, and
Phase~4 answers RQ5. RQ4 is a re-analysis of the Phase~3 runs under a paired comparison
rather than a separate experiment, which is why it costs no additional training.

\subsubsection{Phase 1: Full-Graph Baseline}
\label{sec:method:experiment:phase1}
Train the GNN on full, unreduced AIGs to establish predictive accuracy and baseline hardware overhead (VRAM, training time).

\subsubsection{Phase 2: Offline Reduction Profiling}
\label{sec:method:experiment:phase2}
Apply and measure the offline computational cost and shrinkage achieved by the three
reduction categories: Partitioning, Sparsification, and Coarsening/Summarization. This phase
requires no training, and it is what determines which methods are viable at corpus scale at
all --- a reduction whose offline cost exceeds the training it saves is only worthwhile if
the artifact is reused across many runs.

\subsubsection{Phase 3: Trade-off Analysis}
\label{sec:method:experiment:phase3}
Train models on the newly reduced datasets and compare their performance against the Phase 1 baseline to establish a Pareto front (predictive degradation vs. memory/speed gains).

\subsubsection{Phase 4: Cross-State Generalization}
\label{sec:method:experiment:phase4}
Query the models trained exclusively on reduced graphs using normal, full-sized graphs during inference to evaluate structural generalization capabilities or model graph indistinguishability.

\subsection{Reproducibility}
\label{sec:method:experiment:reproducibility}
All stochastic components are seeded at a fixed value: dataset splitting, the random
partitioning and sparsification methods, the random and hash-based summarization methods,
and model initialisation. The deterministic reduction methods (METIS, level slicing,
AND-gate-only, colour refinement, MFFC coarsening) require no seed and produce identical
output on repeated runs by construction.

Seeding makes a single run repeatable; it does not establish that a result is robust. Each
configuration is trained once, so no run-to-run variance is available, and differences
between configurations are point estimates
(\ref{sec:method:experiment:metrics:stats}).

Reduction artifacts are cached rather than recomputed, which makes results reproducible
without re-running the offline stage. One caveat must be recorded because it has already
caused a real problem: the summarization cache is keyed by method name only, not by the
parameters the method ran with, so changing a parameter and re-running silently reuses the
existing artifacts and produces a corpus of mixed provenance. Until the cache key includes a
parameter hash, any parameter change requires deleting the method's cache by hand.

% TODO: state availability of code, dataset and reduction artifacts, and pin the abc version
% used for dataset generation -- results depend on it and it is not currently recorded.


% Focus on what you add to the existing method. Explain what you will do and why (and how). Do not forget to characterize your research design. There should be a sub-section on the evaluation.
%For DS students, this normally means using manually labelled or ground truth data. For IS students, it is not always needed to have a separate methodology section. You can also integrate the approach with the results in one section. It depends on your type of research what is best fitting.
% Write about your methodology here. Focus on your own contribution. Indicate exactly how you will assess your work in terms of evaluation.
% % It is possible to use a separate section for the Experimental Setup, which then focuses on all settings used in your experiments. It also possible to address the settings in a sub-section under Methodology.
% \section{Equations}
%         We estimate the deformable template parameters~$\theta_t$ and the deformation fields for every data point using maximum likelihood. Letting~$\mathcal{V} = \{\boldsymbol{v}_i\}$ and~$\mathcal{A} = \{a_i\}$,
%         %
%         \begin{align}
%                 \hat{\theta_t}, \hat{\mathcal{V}} &= \arg \max_{\theta_t, \mathcal{V}} \log p_{\theta_t}(\mathcal{V} | \mathcal{X},  \mathcal{A}) \nonumber \\
%                 &= \arg \max_{\theta_t, \mathcal{V}} \log p_{\theta_t}(\mathcal{X} | \mathcal{V}; \mathcal{A}) + \log p(\mathcal{V}),
%                 \label{eq:logpost}
%         \end{align}
%         %
%         where the first term captures the likelihood of the data and deformations, and the second term controls a prior over the deformation fields.

%         \begin{proof}
%                 Awesome proof.
%         \end{proof}

% \section{Long equations in two columns}
%         Note that equations can fill the width pretty quickly when there are two columns.
%         Different mathemtical environments, beyond the basic \verb|\begin{equation}| may be of use. Taking as an example the binomial formula which overflows in Eq. \eqref{binom:eq}:
%         \begin{equation}
%                 \label{binom:eq}
%                 (1+x)^n = \underbrace{(1+x)\times...\times(1+x)}_{n \text{ times}} = \sum_{k=0}^n \binom{n}{k}x^k = \sum_{k=0}^n\frac{n!}{k!(n-k)!}x^k .
%         \end{equation}

%         The \verb|\begin{multline}| environment is an easy although crude fix:
%         \begin{multline}
%                 (1+x)^n = \underbrace{(1+x)\times...\times(1+x)}_{n \text{ times}} \\
%                 = \sum_{k=0}^n \binom{n}{k}x^k \\
%                 = \sum_{k=0}^n\frac{n!}{k!(n-k)!}x^k .
%         \end{multline}
%         \verb|\begin{align}| tends to give the most control over the final result:
%         \begin{align}
%                 (1+x)^n &= \underbrace{(1+x)\times...\times(1+x)}_{n \text{ times}} \nonumber\\
%                         &= \sum_{k=0}^n \binom{n}{k}x^k \nonumber\\
%                         &= \sum_{k=0}^n\frac{n!}{k!(n-k)!}x^k .
%         \end{align}


% \section{Math styles}
%         Different font styles can be used for equations:
%         \begin{itemize}
%                 \item \verb|$a b c A B C 1 2 3$|: $ a b c A B C 1 2 3 $
%                 \item \verb|$\mathbf{a b c A B C 1 2 3}$|: $ \mathbf{a b c A B C 1 2 3} $
%                 \item \verb|$\mathfrak{a b c A B C 1 2 3}$|: $ \mathfrak{a b c A B C 1 2 3} $
%                 \item \verb|$\mathcal{ABC}$|: $ \mathcal{ABC} $
%                 \item \verb|$\mathbb{ABC}$|: $ \mathbb{ABC} $
%                 \item \verb|$\mathsf{ABC}$|: $ \mathsf{ABC} $
%                 \item \verb|$\mathtt{ABC}$|: $ \mathtt{ABC} $
%         \end{itemize}

%         Text and names in equations should be dealt with the \verb|\text| command, for instance:\\
%         \verb|$\mathcal L_{\text{SuperLoss}}$|: $\mathcal L_{\text{SuperLoss}}$ and not $\mathcal L_{SuperLoss}$.